\section{Background}
\label{sec:background}

\subsection{GPU Execution Model and Clock Domains}

A PyTorch forward pass translates a sequence of \texttt{aten::} operator calls through the
dispatcher into one or more CUDA kernel launches.  Each kernel executes on a set of Streaming
Multiprocessors (SMs), organized into thread blocks (CTAs); each CTA contains multiple warps of
32 threads.  GPU hardware exposes performance counters via CUPTI (CUDA Profiling Tools Interface),
which operates in a hardware-domain clock derived from the GPU's free-running oscillator.

\textbf{The clock domain problem.}  Two GPU profiling tools observe the same kernel execution
through different timing mechanisms.  Nsight Systems (\nsys) reads CUPTI activity records that
carry GPU-domain timestamps (\texttt{start\_ns}, \texttt{end\_ns}) from
\texttt{CUPTI\_ACTIVITY\_KIND\_KERNEL}.  Nsight Compute (\ncu) replays each kernel in a separate
profiling pass, injecting performance counter reads between kernel invocations; its internal
timestamps reflect the replay-mode execution clock, not the original capture clock.  Even on the
same physical GPU, the two clocks diverge by tens of microseconds per kernel due to replay overhead
and kernel serialization.  \textbf{Timestamp joins across these two tools are incorrect and
introduce systematic misattribution.}  \sys avoids this entirely (\secref{sec:attribution}).

NVTX (NVIDIA Tools Extension) ranges are CPU-side annotations pushed and popped by the host
process.  PyTorch's \texttt{torch.autograd.profiler.emit\_nvtx()} context manager inserts
\texttt{nvtxRangePushA("aten::linear")} before and \texttt{nvtxRangePop()} after every
\texttt{aten::} dispatch call, labeling the host-side window during which that operator's CUDA
kernels are launched.  These annotations appear in the \nsys SQLite export as
\texttt{NVTX\_EVENTS} rows with CPU-domain timestamps.

\subsection{PyTorch Compilation Stack}

PyTorch 2.x introduced TorchDynamo and TorchInductor \cite{pytorch2024} as a two-stage compilation
pipeline.  Dynamo intercepts Python bytecode, identifies subgraphs of PyTorch operations, and
traces them to an FX graph in the Aten IR (\texttt{torch.ops.aten.*}).  Inductor then lowers this
graph to Triton-generated GPU kernels, applying fusion heuristics such as pointwise-into-reduction
and elementwise-chaining.

Triton \cite{triton2019} generates fused kernels whose names encode constituent operations, \eg
\texttt{triton\_poi\_fused\_relu\_addmm\_0} for a pointwise kernel fusing \texttt{relu} and
\texttt{addmm}.  When NVTX ranges are absent or ambiguous, these names provide a low-confidence
attribution signal (\secref{sec:heuristic}).

\subsection{Roofline Bottleneck Taxonomy}
\label{sec:roofline}

The Roofline model \cite{williams2009roofline} characterizes a kernel's performance ceiling as
$\min(\text{peak FLOP/s},\ \text{AI} \times \text{peak bandwidth})$, where arithmetic intensity
$\text{AI} = \text{FLOPs} / \text{bytes}$.  A kernel below the bandwidth ceiling is
\emph{memory-bound}; above the ridge point it is \emph{compute-bound}.  \sys operationalizes this
taxonomy at the per-operator level using hardware counters:

\begin{itemize}[leftmargin=1.5em]
  \item \textbf{Tensor Core idle} (\texttt{tensor\_core\_active\_pct} $\approx 0$):
    compute-bound path is not engaging hardware-accelerated units; fix via dtype change
    (\secref{sec:cs:gpt2}).
  \item \textbf{Wave starvation} (Waves/SM~$<$~1, i.e., total thread blocks launched
    $<$ SM count): fewer blocks than SMs leaves some processors completely idle
    regardless of per-SM warp occupancy; observable as low \texttt{achieved\_occupancy}
    combined with small grid dimensions; fix via batch padding or kernel fusion
    (\secref{sec:cs:sdpa}).
  \item \textbf{Register pressure} (\texttt{registers\_per\_thread} $> 128$):
    NVIDIA SMs provide 65,536 registers per SM with a maximum of 2,048 resident
    threads~\cite{cudaprogrammingguide}; at 128 registers per thread the
    active-thread ceiling is $\lfloor 65{,}536 / 128 \rfloor = 512$ threads
    (25\% of peak occupancy), severely limiting latency hiding; often resolves
    with dtype reduction (\texttt{bfloat16} halves per-thread register demand on
    GEMM kernels).
  \item \textbf{Cache miss pressure} (low \texttt{l2\_hit\_rate} coincident with
    \texttt{dram\_throughput\_pct}~$>$~60\%): the working set exceeds L2 capacity
    and the kernel is saturating DRAM bandwidth---placing it on the memory-bandwidth
    ceiling of the Roofline model~\cite{williams2009roofline}; on GEMM paths this
    typically signals a column-major vs.\ row-major layout mismatch; fix via weight
    pre-transposition.
  \item \textbf{Memory spill} (\texttt{local\_memory\_spills} $> 0$):
    register-file overflow to L1; fix via kernel substitution (FlashAttention-2 vs.\ standard FMHA).
\end{itemize}

These five patterns cover the bottlenecks observed across the evaluated workloads in
\secref{sec:experiments}.

